\documentclass[noanswers, addpoints, 12pt]{exam}
\usepackage[T1]{fontenc}
\usepackage[utf8]{inputenc}
\usepackage{lmodern}
\usepackage[french]{babel}

\usepackage{graphicx}
\usepackage{tabularx}
\usepackage{minted}
\usepackage{eurosym}
\usepackage{enumitem}
\usepackage{amsmath}
\usepackage{tikz}
\usetikzlibrary{positioning, shapes.multipart, arrows.meta, shadows}

\newcommand{\sql}[1] {\textsc{#1}}

% Configuration TikZ pour le schéma relationnel
\tikzset{
    table/.style={
        draw, rectangle split, rectangle split parts=2,
        rounded corners, fill=gray!5, drop shadow,
        align=left, font=\small\sffamily, minimum width=3.5cm
    }
}

\title{CC Term NSI - SQL : Gestion d'une Bibliothèque}
\author{Babash - NSI Tani Malandi et Johanne Bakalara}
\date{Février 2026, Durée: 1h}

\begin{document}
\runningfooter{}{\footnotesize Babash - NSI Tani Malandi et Johanne Bakalara \textperiodcentered{} CC BY-SA 4.0 \textperiodcentered{} page \thepage}{}
\firstpagefooter{}{\footnotesize Babash - NSI Tani Malandi et Johanne Bakalara \textperiodcentered{} CC BY-SA 4.0 \textperiodcentered{} page \thepage}{}

\maketitle

\pagestyle{headandfoot}
\thispagestyle{headandfoot}
\firstpageheader{Nom : \\  Prénom :\\ Classe :}{}{}

\section*{Consignes:}
\begin{center} 
\fbox{\parbox{\linewidth}{
\begin{center}
Le devoir est à rédiger directement sur le sujet.\\
Toute tentative ou suspicion de triche sera sanctionnée d'un zéro.\\
La calculatrice ou tout autre appareil électronique n'est pas autorisée.\\
Le langage SQL n'est pas sensible à la casse, mais on utilisera les majuscules pour les mots-clés.\\
Chaque instruction doit être terminée par un point-virgule.
\end{center}
}}
\end{center}

\section*{Modèle relationnel :}

\noindent Le système repose sur trois relations (tables) :

\begin{itemize}
    \item \textbf{Auteur} (\underline{id\_auteur}, nom, prenom, nationalite)
    \item \textbf{Livre} (\underline{id\_livre}, titre, genre, prix, \textit{\#id\_auteur})
    \item \textbf{Emprunt} (\underline{id\_emprunt}, date\_emprunt, nom\_emprunteur, \textit{\#id\_livre})
\end{itemize}
\vfill
\begin{center}
\begin{tikzpicture}[node distance=1.5cm and 2cm, >=Stealth, scale=0.9, every node/.style={transform shape}]
    % Table AUTEUR
    \node[table] (auteur) { 
        \textbf{AUTEUR} 
        \nodepart{second} 
        \underline{id\_auteur} : INT\\ 
        nom : TEXT \\ 
        prenom : TEXT \\ 
        nationalite : TEXT 
    };

    % Table LIVRE
    \node[table, right=of auteur] (livre) { 
        \textbf{LIVRE} 
        \nodepart{second} 
        \underline{id\_livre} : INT\\ 
        titre : TEXT \\ 
        genre : TEXT \\ 
        prix : FLOAT \\ 
        \#id\_auteur : INT 
    };

    % Table EMPRUNT
    \node[table, right=of livre] (emprunt) { 
        \textbf{EMPRUNT} 
        \nodepart{second} 
        \underline{id\_emprunt} : INT\\ 
        date\_emprunt : TEXT \\ 
        nom\_emprunteur : TEXT \\ 
        \#id\_livre : INT 
    };

    % Flèches précises des attributs FK vers PK
    % De Livre.#id_auteur (dernière ligne) vers Auteur.id_auteur (première ligne)
    \draw[->, thick, black] ([yshift=-2.9cm]livre.north west) -- ++(-0.6,0) |- ([yshift=-0.8cm]auteur.north east);
    
    % De Emprunt.#id_livre (dernière ligne) vers Livre.id_livre (première ligne)
    \draw[->, thick, black] ([yshift=-2.4cm]emprunt.north west) -- ++(-0.6,0) |- ([yshift=-0.8cm]livre.north east);

\end{tikzpicture}
\end{center}
\vfill
\newpage

\section{Partie 1 : QCM (7 points)}
\textit{Choisir la bonne réponse pour chaque situation.}

\begin{questions}

\question[1] Quelle requête permet d'ajouter l'auteur "Victor Hugo" ?
\begin{choices}
    \choice \begin{minted}{sql}
INSERT INTO Auteur VALUES (1, 'Hugo', 'Victor', 'France');
\end{minted}
    \choice \begin{minted}{sql}
UPDATE Auteur SET nom='Hugo' WHERE prenom='Victor';
\end{minted}
    \choice \begin{minted}{sql}
CREATE Auteur (1, 'Hugo', 'Victor', 'France');
\end{minted}
\end{choices}

\question[1] On souhaite supprimer l'emprunt dont l'id est 102.
\begin{choices}
    \choice \begin{minted}{sql}
DROP FROM Emprunt WHERE id_emprunt = 102;
\end{minted}
    \choice \begin{minted}{sql}
DELETE FROM Emprunt WHERE id_emprunt = 102;
\end{minted}
    \choice \begin{minted}{sql}
UPDATE Emprunt DELETE WHERE id_emprunt = 102;
\end{minted}
\end{choices}

\question[1] Comment supprimer totalement la structure de la table \texttt{Emprunt} ?
\begin{choices}
    \choice \begin{minted}{sql}
DELETE TABLE Emprunt;
\end{minted}
    \choice \begin{minted}{sql}
REMOVE TABLE Emprunt;
\end{minted}
    \choice \begin{minted}{sql}
DROP TABLE Emprunt;
\end{minted}
\end{choices}

\question[1] Quelle requête modifie le prix du livre n°50 pour le passer à 12.50 ?
\begin{choices}
    \choice \begin{minted}{sql}
UPDATE Livre SET prix = 12.50 WHERE id_livre = 50;
\end{minted}
    \choice \begin{minted}{sql}
INSERT INTO Livre (prix) VALUES (12.50) WHERE id_livre = 50;
\end{minted}
    \choice \begin{minted}{sql}
SET prix = 12.50 FROM Livre WHERE id_livre = 50;
\end{minted}
\end{choices}

\question[1] Pour sélectionner les livres dont le genre est 'BD' ou 'Manga' :
\begin{choices}
    \choice \begin{minted}{sql}
WHERE genre IN ('BD', 'Manga');
\end{minted}
    \choice \begin{minted}{sql}
WHERE genre = 'BD' AND genre = 'Manga';
\end{minted}
    \choice \begin{minted}{sql}
WHERE genre LIKE 'BD' 'Manga';
\end{minted}
\end{choices}

\question[1] Quelle fonction permet de calculer la moyenne des prix des livres ?
\begin{choices}
    \choice \begin{minted}{sql}
SUM(prix)
\end{minted}
    \choice \begin{minted}{sql}
AVG(prix)
\end{minted}
    \choice \begin{minted}{sql}
COUNT(prix)
\end{minted}
\end{choices}

\question[1] Comment obtenir seulement les 3 livres les moins chers ?
\begin{choices}
    \choice \begin{minted}{sql}
ORDER BY prix ASC LIMIT 3;
\end{minted}
    \choice \begin{minted}{sql}
ORDER BY prix DESC LIMIT 3;
\end{minted}
    \choice \begin{minted}{sql}
SELECT 3 FROM Livre ORDER BY prix;
\end{minted}
\end{choices}

\newpage

\section{Partie 2 : Rédaction de requêtes (14 points)}
\textit{Dans cette partie, rédigez les requêtes correspondant aux demandes formulées. Des points peuvent être attribués pour des requêtes non-fonctionnelles mais partiellement juste.}

\question[2] Écrire la requête permettant de créer la table \textbf{Emprunt} en respectant le modèle fourni (types et clés primaire et étrangère).
\vspace{4cm}

\question[2] Donner la requête permettant de lister uniquement les titres de tous les livres.
\vfill

\question[2] Donner la requête permettant d'afficher les titres des livres dont le genre est 'Roman' et dont le prix est strictement supérieur à 20 €. 
\vfill

\question[2] Écrire la requête permettant de connaître le nombre total d'enregistrements dans la table \textbf{Emprunt}. 
\vfill

\question[2] Écrire la requête permettant d'afficher le titre du livre et le nom de l'auteur de chaque ouvrage ayant un prix supérieur à 50 €. 
\vspace{3cm}


\question[2] Donner la requête permettant de sélectionner les titres des livres qui contiennent le mot "Secret" dans leur titre.
\vfill

\question[2] Écrire la requête affichant toutes les informations des livres, triés par ordre alphabétique des titres. 
\vfill

\end{questions}
\end{document}